\chapter{Assemblaggio aplotipi da frammenti. Definizione del problema MEC (Minimum Error
Correction) e formulazione con il grafo bipartito}

L'assemblaggio di aplotipi è un problema computazionale che consiste nel ricostruire i due aplotipi di un individuo diploide a partire da read ottenute dal sequenziamento.
Nel caso diploide, ogni individuo possiede due copie di ciascun cromosoma, una di origine materna e una di origine paterna. Le read prodotte dal sequenziamento provengono da entrambi gli aplotipi, ma non si conosce a priori da quale dei due provenga ciascuna read.
\begin{itemize}
    \item \textbf{Input}: una collezione di read/frammenti ottenuti dal sequenziamento e allineati al reference.
    \item \textbf{Output}: una bipartizione delle read in due insiemi, ciascuno corrispondente a uno dei due aplotipi.

\end{itemize}


Dal punto di vista computazionale, l'input viene rappresentato tramite una matrice dei frammenti \(F\). Le righe ($n$)rappresentano i frammenti \(f_1,\dots,f_n\), mentre le colonne($m$) rappresentano le posizioni SNP considerate.
\[
F \in \{1,0,-\}^{n \times m}
\]
Un elemento \(F(i,j)=0\) indica che il frammento \(f_i\) osserva allele 0 nella posizione SNP \(j\), mentre \(F(i,j)=1\) indica che osserva allele 1. Il simbolo \(F(i,j)=-\) indica invece che il frammento non copre quella posizione.
Le posizioni omozigoti possono essere ignorate, perché non aiutano a distinguere i due aplotipi. Sono invece informative le posizioni eterozigoti, dove i due aplotipi hanno valori diversi.
In assenza di errori, l'obiettivo è trovare una bipartizione dei frammenti in due insiemi \(R\) e \(S\), tali che:
\[
R \cup S = \{f_1,\dots,f_n\}, \qquad R \cap S = \emptyset.
\]

I frammenti nello stesso insieme devono essere compatibili, 
cioè non devono presentare valori opposti in una stessa colonna SNP coperta da entrambi.

Due frammenti \(f_i\) e \(f_k\) sono compatibili se, per ogni colonna \(j\),
\[
F(i,j)=- 
\quad \text{oppure} \quad 
F(k,j)=- 
\quad \text{oppure} \quad 
F(i,j)=F(k,j).
\]
Al contrario, due frammenti sono in conflitto se esiste almeno una colonna \(j\) tale che entrambi osservano quella 
posizione e riportano valori opposti:
\[
\exists j :
F(i,j)\neq -, \qquad
F(k,j)\neq -, \qquad
F(i,j)\neq F(k,j).
\]
A partire da questa formulazione è possibile descrivere il problema tramite un grafo. Si costruisce il grafo dei conflitti:
\[
G_F=(V,E),
\]
dove i vertici sono i frammenti della matrice:
\[
V=\{f_1,\dots,f_n\}.
\]
E si inserisce un arco tra due frammenti \(f_i\) e \(f_k\) se e solo se essi sono in conflitto. 
\[
(f_i,f_k)\in E
\iff
\exists j :
F(i,j)\neq -, \quad
F(k,j)\neq -, \quad
F(i,j)\neq F(k,j).
\] 
Quindi un arco del grafo collega due frammenti che non possono stare nello stesso aplotipo.

Il punto fondamentale è che, in assenza di errori, esiste una bipartizione dei frammenti in due insiemi senza conflitti se e solo se il grafo dei conflitti \(G_F\) è bipartito:

\[
\exists (R,S) \text{ bipartizione conflict-free di } F
\iff
G_F \text{ è bipartito}.
\]

Infatti, se esiste una bipartizione corretta, allora non possono esserci archi tra frammenti appartenenti allo stesso insieme, perché un arco rappresenta un conflitto. 
Dunque tutti gli archi collegano frammenti appartenenti a insiemi diversi.
Viceversa, se il grafo dei conflitti è bipartito, allora i suoi vertici possono essere divisi in due insiemi i
ndipendenti \(R\) e \(S\). Poiché non ci sono archi interni a \(R\) né a \(S\), 
i frammenti nello stesso insieme non sono in conflitto e possono quindi essere assegnati allo stesso aplotipo.

La costruzione del grafo dei conflitti è polinomiale. Infatti, il grafo ha tanti vertici quante sono le righe della 
matrice \(F\), quindi \(n\) vertici. Gli archi collegano coppie di frammenti in conflitto; nel caso peggiore 
si confrontano tutte le coppie di frammenti, cioè \(O(n^2)\) coppie, controllando le \(m\) colonne SNP. 
Quindi la costruzione del grafo richiede tempo \(O(n^2m)\). Anche verificare se il grafo è bipartito 
si può fare in tempo polinomiale, quindi in assenza di errori la bipartizione dei frammenti è in \(P\).

Adesso si può formulare la definizone completa del problema dell'assemblaggio di aplotipi:
\begin{itemize}
    \item \textbf{Input}: matrice dei frammenti \(F\) di $m$ colonne (SNP) ed $n$ righe $f_1,f_2 \dots f_n$.
    \item \textbf{Output}: trovare una bipartizione dei frammenti delle righe di \(F\) in due insiemi \(R\) e \(S\) tale che $\forall f_i, fj \in R (S)$ non siano in conflitto.
\end{itemize}

In presenza di errori di sequenziamento, però, una tale bipartizione \(R,S\) può non esistere. 
Infatti una read può contenere un valore errato in una posizione SNP, cioè può riportare \(0\) al posto di \(1\), 
oppure \(1\) al posto di \(0\). Questo può introdurre conflitti artificiali nella matrice dei frammenti e impedire 
l'esistenza di una bipartizione senza conflitti. Per questo viene introdotto il modello \textbf{MEC}, 
cioè \emph{Minimum Error Correction}.
\\
Il problema MEC consiste nel trovare le posizioni \((i,j)\) della matrice \(F\) da correggere in modo che, d
opo tali correzioni, esista una bipartizione dei frammenti in due insiemi \(R\) e \(S\) senza conflitti. 
Una correzione consiste nello \emph{switch} di un valore osservato:
\[
0 \leftrightarrow 1.
\]
Indicando con \(P\) l'insieme delle posizioni corrette, l'obiettivo del problema MEC è minimizzare il numero di tali posizioni:
\[
\min |P|.
\]
Quindi MEC cerca il minimo numero di valori della matrice da cambiare affinché sia possibile ottenere la 
bipartizione dei frammenti nei due aplotipi.

Se invece la matrice \(F\) è pesata, cioè se a ogni entry osservata \(F(i,j)\) è associato un peso
\[
w(i,j),
\]
si ottiene la versione pesata del problema, detta \textbf{WMEC}, cioè \emph{Weighted Minimum Error Correction}. 
In questo caso il peso \(w(i,j)\) rappresenta il costo di correggere la posizione \((i,j)\). 
L'obiettivo non è più minimizzare semplicemente il numero di correzioni, ma minimizzare la somma dei pesi delle 
posizioni corrette. Se \(P\) è l'insieme delle posizioni di switch, allora l'obiettivo diventa:
\[
\min \sum_{(i,j)\in P} w(i,j).
\]
Di conseguenza, MEC minimizza il numero di correzioni, mentre WMEC minimizza il costo totale delle correzioni. 
Il problema WMEC è computazionalmente difficile: è NP-difficile, e la sua versione decisionale è NP-completa.

Per affrontare WMEC si può usare una programmazione dinamica sulle colonne della matrice. 
L'idea è costruire progressivamente la soluzione considerando prima le colonne \(1,\dots,j\), 
poi estendendo la soluzione alla colonna \(j+1\). Una soluzione ottima sulle prime \(j+1\) 
colonne contiene una soluzione ottima sulle prime \(j\) colonne; per questo si memorizza il costo minimo 
delle possibili sottomatrici già considerate.
Sia
\[
B=(R,S)
\]
una bipartizione dei frammenti attivi nella colonna \(j\), cioè dei frammenti che in quella colonna non hanno il 
simbolo ``\(-\)''. Una bipartizione successiva
\[
(R',S')
\]
estende
\[
(R,S)
\]
se conserva le assegnazioni già fatte e può eventualmente aggiungere nuovi frammenti. Formalmente:
\[
R' \supseteq R,
\qquad
S' \supseteq S.
\]

Si introduce allora il costo
\[
C(j,B),
\]
dove \(B=(R,S)\). Questo valore rappresenta il costo minimo per correggere la sottomatrice formata dalle colonne 
\(1,\dots,j\), in modo che sia senza conflitti e sia coerente con la bipartizione \(B\) nella colonna \(j\). 
Resta quindi da definire il costo locale di correzione di una colonna. 
Fissata una bipartizione \((R,S)\), si indica con
\[
\Delta(j,(R,S))
\]
il costo necessario per rendere senza conflitti la colonna \(j\) rispetto alla bipartizione \((R,S)\).
Per calcolare questo costo si definisce, per \(x\in\{0,1\}\),
\[
W^x(j,R)=
\sum_{\substack{i\in R\\F(i,j)=x}} w(i,j).
\]
Questa quantità rappresenta la somma dei pesi delle posizioni nella colonna \(j\), 
limitatamente ai frammenti in \(R\), che hanno valore \(x\). Analogamente si definisce \(W^x(j,S)\) 
per i frammenti dell'insieme \(S\).

Senza imporre l'ipotesi che la colonna sia eterozigote, per rendere conflict-free la colonna \(j\) si può scegliere 
indipendentemente se portare i frammenti di \(R\) al valore \(0\) oppure al valore \(1\), e 
lo stesso per \(S\). Per questo il costo locale è:
\[
\Delta(j,(R,S))
=
\min\{W^0(j,R),W^1(j,R)\}
+
\min\{W^0(j,S),W^1(j,S)\}.
\]
Questa formula sceglie, per ciascun insieme della bipartizione, il valore che richiede il costo minore di correzione.
Nel caso \emph{all heterozygous}, invece, si assume che le colonne considerate siano eterozigoti. 
Quindi i due insiemi della bipartizione devono rappresentare alleli diversi. 
Se \(R\) viene portato a \(0\), allora \(S\) deve essere portato a \(1\); se \(R\) viene portato a \(1\), 
allora \(S\) deve essere portato a \(0\). Le due possibilità sono:
\[
R \to 0,\quad S \to 1,
\]
oppure:
\[
R \to 1,\quad S \to 0.
\]
Con questa assunzione, il costo locale diventa:
\[
\Delta(j,(R,S))
=
\min
\left\{
W^0(j,R)+W^1(j,S),
\;
W^1(j,R)+W^0(j,S)
\right\}.
\]
Questa formula impone che i due gruppi della bipartizione abbiano valori opposti nella colonna \(j\), come richiesto da una posizione eterozigote.
La ricorrenza della programmazione dinamica si basa su questa idea: per calcolare il costo fino alla colonna \(j+1\), si somma il costo locale della colonna \(j+1\) al miglior costo ottenuto fino alla colonna \(j\), considerando solo le bipartizioni compatibili. Data una bipartizione \((R,S)\) per la colonna \(F(j+1)\), si ha:
\[
C(j+1,(R,S))
=
\Delta(j+1,(R,S))
+
\min_B C(j,B),
\]
dove \(B\) varia tra le bipartizioni di interesse della colonna \(j\) compatibili con \((R,S)\) sulla colonna \(j+1\).
La compatibilità serve a garantire che i frammenti comuni tra due colonne consecutive non cambino insieme. Se un frammento compare sia nella colonna \(j\) sia nella colonna \(j+1\), e nella colonna \(j\) era assegnato a \(R\), allora deve rimanere nello stesso lato anche nella colonna successiva. Lo stesso vale per i frammenti assegnati a \(S\).
In questo modo, WMEC viene risolto costruendo progressivamente soluzioni ottime per prefissi della matrice, mantenendo la coerenza delle bipartizioni tra colonne consecutive e minimizzando il costo totale delle correzioni.